\documentclass[12pt,a4paper]{article}
\usepackage{graphicx} % Required for inserting images
\usepackage{chemfig}
\usepackage{float}
\usepackage{lewis}
\author{Svein-Tore Narvestad}
\date{September 2025}

\begin{document}
\textbf{\large{Prøve om kjemiske bindinger:}}\\\\\\
1.	Hva menes med begrepet elektronegativitet og hvor finner vi grunnstoffene som har høy/lav elektronegativitet?\\\\
.\\
2.	Hva menes med ioniseringsenergi? Hvilke grunnstoff har lav\\ ioniseringsenergi?\\\\
\noindent
3. Forklar hvorfor karbondioksid ($CO_2$) ikke er en dipol.\\\\
\noindent
4.	Hva menes med hydrogenbinding?  Nevn 3 stoff som har hydrogenbinding.\\\\
\noindent
5.	 Gitt følgende stoffer: $CaCl_2, KCl, CO_2,CH_4$ og NaCl.
Hva er elektronegativitetsverdien for grunnstoffene i stoffene over?\\\\
\noindent
6.	Hvilken bindingstype har vi mellom molekylene i stoffene $CaCl_2, KCl, CO_2,CH_4$ og NaCl?\\\\

\noindent
7.	Tegn elektronprikkstrukturen for stoffene $H_2O, NH_3, CO_2, CH_4$ og NaCl.\\\\
\noindent
8.	Er noen av molekylene dipoler? Begrunn svaret.\\\\
\begin{figure}[H]
    \centering
    \includegraphics[width=0.5\linewidth]{elnegverdier.png}
    
\end{figure}
\newpage
\textbf{Oversikt over type binding baser på forskjell i elektronegativitet:}
\begin{figure}[H]
    \centering
    \includegraphics[width=0.8\linewidth]{image.png}
    
    
\end{figure}
\end{document} 
